\subsection{\foreign{Pseudo Boolean}}

Pada akhirnya, sampailah di penghujung desain algoritma dimensi metrik menggunakan SMT. Cermati bentuk \foreign{constraint} yang terdapat dalam \foreign{linear integer arithmetic} sudah sesuai dengan representasi \foreign{constraint} yang dibutuhkan dalam \foreign{pseudo boolean}. Maka dari itu, konstruksi \foreign{constraint} dapat diturunkan kembali menggunakan variabel boolean dari pada \foreign{integer}, tetapi menggunakan teori yang lebih tinggi dari \foreign{boolean satisfiability}, yaitu \foreign{pseudo boolean}. Di sini, didapatkan kembali representasi natural keanggotaan sebagai nilai \(\True\), dan secara bersamaan dapat diekspresikan bentuk \foreign{constraint} yang memuat logika kardinalitas secara lebih natural. Dengan demikian, \autoref{eq:linear-integer-arithmetic-cardinality} dan \autoref{eq:linear-integer-arithmetic} dapat translasikan langsung menjadi pseudo boolean sebagai berikut:

\begin{equation}
	\sum \limits_{\displaystyle v \in V} \boolmap(v) \enspace = k
	\label{eq:pseudo-boolean-cardinality}
\end{equation}

\begin{equation}
	\bigwedge \limits_{\displaystyle P_{*j} \in P} \enspace \sum \limits_{\displaystyle v \in V[P_{*j}]} \boolmap(v) \enspace \leq k - 1
	\label{eq:pseudo-boolean}
\end{equation}

Meskipun perubahan yang dilakukan hanyalah membawa turun dari \foreign{linear integer arithmetic} ke \foreign{pseudo boolean}, dengan bentuk \foreign{constraint} yang kurang lebih sama persis, ketika tidak ada \foreign{expresiveness} yang hilang, maka secara umum \foreign{solver} yang menggunakan \foreign{pseudo boolean} akan lebih efektif dibandingkan yang menggunakan \foreign{linear integer arithmetic}.

\begin{example}
	Ambil graf \autoref{fig:algorithm-example}. Jika \(\boolmap(v_i) = x_i\), maka \foreign{constraint} yang berkorespondensi dengan \foreign{resolving set} berkardinalitas \(k\) adalah:
	\[
		\smash[b]{
		\begin{array}{c}
			x_1 + x_2 + x_3 + x_4 + x_5 = k \\ [8pt]
			x_3 + x_4 + x_5 \leq k - 1 \\
			x_2 + x_4 + x_5 \leq k - 1 \\
			x_1 + x_4 + x_5 \leq k - 1 \\
			x_1 + x_2 + x_3 \leq k - 1
		\end{array}
		} \vspace{.6em}
	\]
	\label{ex:pseudo-boolean-constraint}
\end{example}

Sehingga, didapat desain algoritma yang menggunakan \foreign{pseudo boolean}.

\begin{example}
	Untuk graf \autoref{fig:algorithm-example}, mencari dimensi metriknya menggunakan \autoref{alg:pseudo-boolean}, dapat dilakukan dengan mengambil pada setiap tahapannya \foreign{constraint} yang telah tertera pada \autoref{ex:pseudo-boolean-constraint}.
	\begin{itemize}
		\item Untuk \(k = 5\), ditemukan \(\sat\).
		\item Untuk \(k = 4\), ditemukan \(\sat\).
		\item Untuk \(k = 3\), ditemukan \(\sat\).
		\item Untuk \(k = 2\), ditemukan \(\unsat\).
	\end{itemize}
	Disimpulkan bahwa \(2 + 1 = 3\) adalah dimensi metrik dari graf tersebut.
\end{example}

\begin{algorithm}
	\foreign{Input}:
	\begin{itemize}
		\item Representasi graf dalam representasi matriks ketetanggaan \(M\).
		\item Mencari dimensi metrik atau dimensi metrik sisi?
	\end{itemize}
	Prosedur:
	\begin{enumerate}
		\item[1--5.] Sama seperti pada \autoref{alg:linear-integer-arithmetic}.
		\setcounter{enumi}{5}
		\item Mulai variabel \(k\) dari banyaknya simpul, yaitu \(k := |V|\).
		\item Selama \(k \geq 0\):
			\begin{enumerate}
					\item Tambahkan \foreign{constraint} berdasarkan \autoref{eq:pseudo-boolean-cardinality}, dan berdasarkan \autoref{eq:pseudo-boolean} dari \(P\), ke \foreign{solver} SMT.
					\item Jalankan \foreign{solver} untuk mencari \foreign{model}.
					\item Jika dihasilkan \(\unsat\), lompat ke tahap (9).
					\item Hapus \foreign{constraint} yang ditambahkan dari tahap (7a) sebelumnya.
					\item Kurangi \(k\) sebanyak satu, yaitu \(k := k - 1\).
			\end{enumerate}
		\item Simpulkan bahwa 0 adalah dimensi metrik atau dimensi metrik sisi, dan \texttt{BERHENTI}.
		\item Simpulkan bahwa \(k + 1\) adalah dimensi metrik atau dimensi metrik sisi.
	\end{enumerate}
	\caption{\foreign{Pseudo Boolean}}
	\label{alg:pseudo-boolean}
\end{algorithm}
